% BHU PYQ -- Professional Exam Template
\documentclass[11pt,a4paper]{article}

\usepackage[a4paper,top=2.5cm,bottom=2.5cm,left=2.8cm,right=2.8cm,headheight=14pt]{geometry}
\usepackage{amsmath,amssymb}
\usepackage[T1]{fontenc}
\usepackage[utf8]{inputenc}
\usepackage{lmodern}
\usepackage{microtype}
\usepackage{fancyhdr}
\usepackage{enumitem}
\usepackage{listings}
\usepackage{hyperref}
\usepackage{lastpage}
\usepackage{parskip}

\hypersetup{colorlinks=true,urlcolor=black,linkcolor=black,
  pdftitle={BVCA-304: Java Applet Programming -- BHU PYQ}}

\pagestyle{fancy}
\fancyhf{}
\renewcommand{\headrulewidth}{0.4pt}
\renewcommand{\footrulewidth}{0.4pt}
\fancyhead[L]{\small\textit{Banaras Hindu University}}
\fancyhead[R]{\small\textit{BVCA-304: Java Applet Programming}}
\fancyfoot[C]{\small Page \thepage\ of \pageref{LastPage}}

\lstset{
  basicstyle=\small\ttfamily,
  frame=single,
  breaklines=true,
  columns=flexible,
  keepspaces=true
}

\newlist{parts}{enumerate}{2}
\setlist[parts,1]{label=(\alph*),leftmargin=*,itemsep=4pt,topsep=3pt}
\setlist[parts,2]{label=(\roman*),leftmargin=*,itemsep=2pt,topsep=2pt}
\newcommand{\pts}[1]{\hfill\textit{[#1]}}

\begin{document}

\begin{center}
  {\large\bfseries BANARAS HINDU UNIVERSITY}\\[2pt]
  {\normalsize B.Voc. Semester III Examination, 2023-24}\\[6pt]
  \rule{0.8\linewidth}{0.5pt}\\[6pt]
  {\large\bfseries Paper: BVCA-304 --- Java Applet Programming}\\[4pt]
  {\normalsize Time: Three Hours \hfill Maximum Marks: 70}
\end{center}

\rule{\linewidth}{0.4pt}

\smallskip
\noindent\textit{Instructions: Attempt all questions. Symbols have their usual meanings.}

\bigskip

\noindent{\large\bfseries SECTION --- A}\\[4pt]
\rule{\linewidth}{0.4pt}\smallskip

Long Answer type questions to be answered in about 500 words each. Attempt any two: \hfill \textbf{12 $\times$ 2 = 24 Marks}

\medskip
\noindent\textbf{Question 1.} What do you understand by the concept of Object Oriented Language? Explain how Java is specified as an Object Oriented Language.

\medskip
\noindent\textbf{Question 2.} Explain about \texttt{for}, \texttt{while} and \texttt{do-while} loop structures in Java with examples. Why are they categorized as entry-controlled and exit-controlled loops?

\medskip
\noindent\textbf{Question 3.} What are operators? Explain all the Operators supported by the Java programming language.

\medskip
\noindent\textbf{Question 4.} Discuss the Switch Statement. Write a program to calculate the percentage of total marks obtained in 5 subjects by a student and display the equivalent grade as S, A, B, C, D, or F. (Grade criteria: 91--100\% as S grade; 81--90\% as A grade; 71--80\% as B grade and so on. F grade represents Fail).

\bigskip\noindent{\large\bfseries SECTION --- B}\\[4pt]
\rule{\linewidth}{0.4pt}\smallskip

\noindent\textbf{Question 5.} Attempt any six parts (Answer in about 200 words each): \hfill \textbf{6 $\times$ 6 = 36 Marks}
\begin{parts}
  \item What do you understand by classes and objects? How are objects created from a class?
  \item What are arrays? How will you implement one-dimensional and two-dimensional arrays?
  \item Define JDK and JVM. Why is Java considered a platform-independent language?
  \item Differentiate between static, instance, and local variables.
  \item Discuss break and continue statements in Java.
  \item Write a program to count the digits of a number input by the user.
  \item Write a Java program to find the area and circumference of a circle using the concept of objects and classes.
  \item Discuss any four inbuilt methods for String manipulation given in the Java library.
\end{parts}

\bigskip\noindent{\large\bfseries SECTION --- C}\\[4pt]
\rule{\linewidth}{0.4pt}\smallskip

\noindent\textbf{Question 6.} Attempt all parts: \hfill \textbf{10 $\times$ 1 = 10 Marks}
\begin{parts}
  \item What does an Eclipse IDE allow?
  \begin{enumerate}[label=(\alph*),itemsep=2pt]
    \item Editing the program
    \item Compiling the program
    \item Both Editing and Compiling the program
    \item Editing, Building, Debugging and Compiling the program
  \end{enumerate}
  
  \item Which keyword is used for accessing the features of a package?
  \begin{enumerate}[label=(\alph*),itemsep=2pt]
    \item package
    \item import
    \item extends
    \item export
  \end{enumerate}
  
  \item Which of these values is returned by the \texttt{read()} method if end of file (EOF) is encountered?
  \begin{enumerate}[label=(\alph*),itemsep=2pt]
    \item 0
    \item 1
    \item -1
    \item Null
  \end{enumerate}
  
  \item Which one of the following is not a Java feature?
  \begin{enumerate}[label=(\alph*),itemsep=2pt]
    \item Object-oriented
    \item Use of pointers
    \item Portable
    \item Dynamic and Extensible
  \end{enumerate}
  
  \item JRE stands for \_\_\_\_\_\_.
  \begin{enumerate}[label=(\alph*),itemsep=2pt]
    \item Java Run Ecosystem
    \item JDK Runtime Environment
    \item Java Runtime Environment
    \item None of these
  \end{enumerate}
  
  \item Which class in Java is used to take input from the user?
  \begin{enumerate}[label=(\alph*),itemsep=2pt]
    \item Scanner
    \item Input
    \item Applier
    \item None of these
  \end{enumerate}
  
  \item Automatic type conversion is possible in which of the following cases?
  \begin{enumerate}[label=(\alph*),itemsep=2pt]
    \item Byte to int
    \item Int to long
    \item Long to int
    \item Short to int
  \end{enumerate}
  
  \item Identify the output of the following program:
\begin{lstlisting}[language=Java]
String str = "abcde";
System.out.println(str.substring(1, 3));
\end{lstlisting}
  \begin{enumerate}[label=(\alph*),itemsep=2pt]
    \item abc
    \item bc
    \item bcd
    \item cd
  \end{enumerate}
  
  \item Find the output of the following code:
\begin{lstlisting}[language=Java]
int ++a = 100;
System.out.println(++a);
\end{lstlisting}
  \begin{enumerate}[label=(\alph*),itemsep=2pt]
    \item 101
    \item Compile error as ++a is not a valid identifier
    \item 100
    \item None of these
  \end{enumerate}
  
  \item Find the output of the following code:
\begin{lstlisting}[language=Java]
public class Solution {
    public static void main(String args[]) {
        int i;
        for (i = 1; i < 6; i++) {
            if (i > 3) continue;
        }
        System.out.println(i);
    }
}
\end{lstlisting}
  \begin{enumerate}[label=(\alph*),itemsep=2pt]
    \item 3
    \item 4
    \item 5
    \item 6
  \end{enumerate}
\end{parts}

\vfill\vfill
\rule{\linewidth}{0.4pt}
\begin{center}\small
\href{https://bhu-exam.vercel.app/}{Click here to visit our website}\\[4pt]\href{https://me-aryan.vercel.app/}{Click here to visit our contributor}\\[4pt]\href{https://quashan-me.vercel.app/}{Click here to visit our contributor}
\end{center}

\end{document}
